\chapter{Introduction}
\label{ch:introduction}

\section{Context}

The global energy transition toward renewable sources has accelerated in recent years, with solar photovoltaic (PV) installations reaching over 1.6~TW of cumulative capacity worldwide by 2024 \cite{irena2023}. Morocco, through its National Energy Strategy 2030, has committed to achieving 52\% of installed electricity capacity from renewable sources \cite{morocco_energy}. The country's solar potential, particularly in cities such as Meknès, Fès, and Marrakech, makes rooftop PV a viable option for residential and commercial buildings.

However, the integration of distributed PV generation into buildings introduces challenges. Solar production is intermittent and does not naturally align with consumption patterns: PV generation peaks at midday while residential demand peaks in the evening. Without intelligent management, a significant portion of PV energy is either curtailed or exported at unfavorable rates, reducing the economic value of the installation.

Battery energy storage systems (BESS) partially address this mismatch by storing surplus PV energy for later use. However, the effectiveness of a BESS depends heavily on its dispatch strategy. A simple rule-based approach---charge when there is surplus, discharge when there is deficit---fails to account for future consumption and production patterns, leading to suboptimal battery utilization.

\section{Motivation}

This project is motivated by the observation that the combination of three technologies---machine learning forecasting, mathematical optimization, and IoT-based real-time monitoring---can significantly improve the performance of a building energy management system compared to conventional reactive strategies.

Specifically:
\begin{enumerate}[nosep]
    \item \textbf{Forecasting} enables proactive rather than reactive energy dispatch.
    \item \textbf{Optimization} finds the mathematically best schedule given forecasts and constraints.
    \item \textbf{IoT monitoring} provides the real-time data foundation for both.
\end{enumerate}

The academic motivation is to integrate knowledge from renewable energy engineering, electrical engineering, embedded systems, database design, machine learning, and optimization into a single coherent engineering project.

\section{Objectives}

The primary objectives of this project are:

\begin{enumerate}
    \item Design and implement a real-time energy monitoring system using ESP32 and appropriate sensors.
    \item Build a data pipeline that stores, processes, and visualizes energy measurements.
    \item Develop machine learning models that forecast building energy consumption and PV production.
    \item Implement an anomaly detection system for identifying unusual energy patterns.
    \item Design and implement an optimization-based energy management system using MILP.
    \item Quantitatively compare the optimized EMS against a conventional rule-based baseline.
    \item Present the complete system on a professional dashboard.
\end{enumerate}

\section{Report Structure}

This report is organized as follows. Chapter~\ref{ch:problem} defines the problem statement. Chapter~\ref{ch:literature} reviews relevant literature. Chapter~\ref{ch:architecture} presents the system architecture. Chapters~\ref{ch:pv}--\ref{ch:data} detail the PV system, energy management model, IoT architecture, and data acquisition. Chapters~\ref{ch:ml}--\ref{ch:optimization} describe the machine learning methodology, forecasting, anomaly detection, and optimization. Chapter~\ref{ch:simulation} presents the simulation environment. Chapter~\ref{ch:implementation} covers the software implementation. Chapter~\ref{ch:results} presents and analyzes the experimental results. Chapters~\ref{ch:economic} and \ref{ch:environmental} provide economic and environmental analyses. Chapter~\ref{ch:limitations} discusses limitations, and Chapter~\ref{ch:conclusion} concludes with future work.
